\documentclass[11pt]{article}

\title{Math 603: Homework 2}
\author{Swayam Chube}
\date{Last Updated: \today}

\usepackage[utf8]{inputenc} % allow utf-8 input
\usepackage[T1]{fontenc}    % use 8-bit T1 fonts
\usepackage{hyperref}       % hyperlinks
\usepackage{url}            % simple URL typesetting
\usepackage{booktabs}       % professional-quality tables
\usepackage{amsfonts}       % blackboard math symbols
\usepackage{nicefrac}       % compact symbols for 1/2, etc.
\usepackage{microtype}      % microtypography
\usepackage{graphicx}
\usepackage{natbib}
\usepackage{doi}
\usepackage{amssymb}
\usepackage{bbm}
\usepackage{amsthm}
\usepackage{amsmath}
\usepackage{xcolor}
\usepackage{theoremref}
\usepackage{enumitem}
\usepackage{fouriernc}
\usepackage{mdframed}
\usepackage{mathrsfs}
\setlength{\marginparwidth}{2cm}
\usepackage{todonotes}
\usepackage{stmaryrd}
\usepackage[all,cmtip]{xy} % For diagrams, praise the Freyd-Mitchell theorem 
\usepackage{marvosym}
\usepackage{geometry}
\usepackage{titlesec}
\usepackage{mathtools}
\usepackage{tikz}
\usetikzlibrary{cd}
\usepackage{epigraph}
\setlength\epigraphwidth{0.4\textwidth}

\renewcommand{\qedsymbol}{$\blacksquare$}
% \renewcommand{\familydefault}{\sfdefault} % Do you want this font? 

% Uncomment to override  the `A preprint' in the header
% \renewcommand{\headeright}{}
% \renewcommand{\undertitle}{}
% \renewcommand{\shorttitle}{}

\hypersetup{
    pdfauthor={Swayam Chube},
    colorlinks=true,
	citecolor=blue,
}

\newtheoremstyle{thmstyle}%               % Name
  {}%                                     % Space above
  {}%                                     % Space below
  {}%                             % Body font
  {}%                                     % Indent amount
  {\bfseries\scshape}%                            % Theorem head font
  {.}%                                    % Punctuation after theorem head
  { }%                                    % Space after theorem head, ' ', or \newline
  {\thmname{#1}\thmnumber{ #2}\thmnote{ (#3)}}%                                     % Theorem head spec (can be left empty, meaning `normal')

\newtheoremstyle{defstyle}%               % Name
  {}%                                     % Space above
  {}%                                     % Space below
  {}%                                     % Body font
  {}%                                     % Indent amount
  {\bfseries\scshape}%                            % Theorem head font
  {.}%                                    % Punctuation after theorem head
  { }%                                    % Space after theorem head, ' ', or \newline
  {\thmname{#1}\thmnumber{ #2}\thmnote{ (#3)}}%                                     % Theorem head spec (can be left empty, meaning `normal')

\theoremstyle{thmstyle}
\newtheorem{theorem}{Theorem}[section]
\newtheorem{lemma}[theorem]{Lemma}
\newtheorem{proposition}[theorem]{Proposition}

\theoremstyle{defstyle}
\newtheorem{definition}[theorem]{Definition}
\newtheorem{corollary}[theorem]{Corollary}
\newtheorem{porism}[theorem]{Porism}
\newtheorem{remark}[theorem]{Remark}
\newtheorem{interlude}[theorem]{Interlude}
\newtheorem{example}[theorem]{Example}
\newtheorem*{notation}{Notation}
\newtheorem*{claim}{Claim}

% Common Algebraic Structures
\newcommand{\R}{\mathbb{R}}
\newcommand{\Q}{\mathbb{Q}}
\newcommand{\Z}{\mathbb{Z}}
\newcommand{\N}{\mathbb{N}}
\newcommand{\bbC}{\mathbb{C}} 
\newcommand{\K}{\mathbb{K}} % Base field which is either \R or \bbC
\newcommand{\F}{\mathbb{F}} % Base field which is either \R or \bbC
\newcommand{\calA}{\mathcal{A}} % Banach Algebras
\newcommand{\calB}{\mathcal{B}} % Banach Algebras
\newcommand{\calI}{\mathcal{I}} % ideal in a Banach algebra
\newcommand{\calJ}{\mathcal{J}} % ideal in a Banach algebra
\newcommand{\frakM}{\mathfrak{M}} % sigma-algebra
\newcommand{\calO}{\mathcal{O}} % Ring of integers
\newcommand{\bbA}{\mathbb{A}} % Adele (or ring thereof)
\newcommand{\bbI}{\mathbb{I}} % Idele (or group thereof)

% Categories
\newcommand{\catTopp}{\mathbf{Top}_*}
\newcommand{\catGrp}{\mathbf{Grp}}
\newcommand{\catTopGrp}{\mathbf{TopGrp}}
\newcommand{\catSet}{\mathbf{Set}}
\newcommand{\catTop}{\mathbf{Top}}
\newcommand{\catRing}{\mathbf{Ring}}
\newcommand{\catCRing}{\mathbf{CRing}} % comm. rings
\newcommand{\catMod}{\mathbf{Mod}}
\newcommand{\catMon}{\mathbf{Mon}}
\newcommand{\catMan}{\mathbf{Man}} % manifolds
\newcommand{\catDiff}{\mathbf{Diff}} % smooth manifolds
\newcommand{\catAlg}{\mathbf{Alg}}
\newcommand{\catRep}{\mathbf{Rep}} % representations 
\newcommand{\catVec}{\mathbf{Vec}}

% Group and Representation Theory
\newcommand{\chr}{\operatorname{char}}
\newcommand{\Aut}{\operatorname{Aut}}
\newcommand{\GL}{\operatorname{GL}}
\newcommand{\im}{\operatorname{im}}
\newcommand{\tr}{\operatorname{tr}}
\newcommand{\id}{\mathbf{id}}
\newcommand{\cl}{\mathbf{cl}}
\newcommand{\Gal}{\operatorname{Gal}}
\newcommand{\Tr}{\operatorname{Tr}}
\newcommand{\sgn}{\operatorname{sgn}}
\newcommand{\Sym}{\operatorname{Sym}}
\newcommand{\Alt}{\operatorname{Alt}}

% Commutative and Homological Algebra
\newcommand{\spec}{\operatorname{spec}}
\newcommand{\mspec}{\operatorname{m-spec}}
\newcommand{\Spec}{\operatorname{Spec}}
\newcommand{\MaxSpec}{\operatorname{MaxSpec}}
\newcommand{\Tor}{\operatorname{Tor}}
\newcommand{\tor}{\operatorname{tor}}
\newcommand{\Ann}{\operatorname{Ann}}
\newcommand{\Supp}{\operatorname{Supp}}
\newcommand{\Hom}{\operatorname{Hom}}
\newcommand{\End}{\operatorname{End}}
\newcommand{\coker}{\operatorname{coker}}
\newcommand{\limit}{\varprojlim}
\newcommand{\colimit}{%
  \mathop{\mathpalette\colimit@{\rightarrowfill@\textstyle}}\nmlimits@
}
\makeatother


\newcommand{\fraka}{\mathfrak{a}} % ideal
\newcommand{\frakb}{\mathfrak{b}} % ideal
\newcommand{\frakc}{\mathfrak{c}} % ideal
\newcommand{\frakf}{\mathfrak{f}} % face map
\newcommand{\frakg}{\mathfrak{g}}
\newcommand{\frakh}{\mathfrak{h}}
\newcommand{\frakm}{\mathfrak{m}} % maximal ideal
\newcommand{\frakn}{\mathfrak{n}} % naximal ideal
\newcommand{\frakp}{\mathfrak{p}} % prime ideal
\newcommand{\frakq}{\mathfrak{q}} % qrime ideal
\newcommand{\fraks}{\mathfrak{s}}
\newcommand{\frakt}{\mathfrak{t}}
\newcommand{\frakz}{\mathfrak{z}}
\newcommand{\frakA}{\mathfrak{A}}
\newcommand{\frakB}{\mathfrak{B}}
\newcommand{\frakI}{\mathfrak{I}}
\newcommand{\frakJ}{\mathfrak{J}}
\newcommand{\frakK}{\mathfrak{K}}
\newcommand{\frakL}{\mathfrak{L}}
\newcommand{\frakN}{\mathfrak{N}} % nilradical 
\newcommand{\frakO}{\mathfrak{O}} % dedekind domain
\newcommand{\frakP}{\mathfrak{P}} % Prime ideal above
\newcommand{\frakQ}{\mathfrak{Q}} % Qrime ideal above 
\newcommand{\frakR}{\mathfrak{R}} % jacobson radical
\newcommand{\frakU}{\mathfrak{U}}
\newcommand{\frakV}{\mathfrak{V}}
\newcommand{\frakW}{\mathfrak{W}}
\newcommand{\frakX}{\mathfrak{X}}

% General/Differential/Algebraic Topology 
\newcommand{\scrA}{\mathscr{A}}
\newcommand{\scrB}{\mathscr{B}}
\newcommand{\scrF}{\mathscr{F}}
\newcommand{\scrM}{\mathscr{M}}
\newcommand{\scrN}{\mathscr{N}}
\newcommand{\scrP}{\mathscr{P}}
\newcommand{\scrO}{\mathscr{O}} % sheaf
\newcommand{\scrR}{\mathscr{R}}
\newcommand{\scrS}{\mathscr{S}}
\newcommand{\bbH}{\mathbb H}
\newcommand{\Int}{\operatorname{Int}}
\newcommand{\psimeq}{\simeq_p}
\newcommand{\wt}[1]{\widetilde{#1}}
\newcommand{\RP}{\mathbb{R}\text{P}}
\newcommand{\CP}{\mathbb{C}\text{P}}

% Miscellaneous
\newcommand{\wh}[1]{\widehat{#1}}
\newcommand{\calM}{\mathcal{M}}
\newcommand{\calP}{\mathcal{P}}
\newcommand{\onto}{\twoheadrightarrow}
\newcommand{\into}{\hookrightarrow}
\newcommand{\Gr}{\operatorname{Gr}}
\newcommand{\Span}{\operatorname{Span}}
\newcommand{\ev}{\operatorname{ev}}
\newcommand{\weakto}{\stackrel{w}{\longrightarrow}}

\newcommand{\define}[1]{\textcolor{blue}{\textit{#1}}}
% \newcommand{\caution}[1]{\textcolor{red}{\textit{#1}}}
\newcommand{\important}[1]{\textcolor{red}{\textit{#1}}}
\renewcommand{\mod}{~\mathrm{mod}~}
\renewcommand{\le}{\leqslant}
\renewcommand{\leq}{\leqslant}
\renewcommand{\ge}{\geqslant}
\renewcommand{\geq}{\geqslant}
\newcommand{\Res}{\operatorname{Res}}
\newcommand{\floor}[1]{\left\lfloor #1\right\rfloor}
\newcommand{\ceil}[1]{\left\lceil #1\right\rceil}
\newcommand{\gl}{\mathfrak{gl}}
\newcommand{\ad}{\operatorname{ad}}
\newcommand{\Stab}{\operatorname{Stab}}
\newcommand{\bfX}{\mathbf{X}}
\newcommand{\Ind}{\operatorname{Ind}}
\newcommand{\bfG}{\mathbf{G}}
\newcommand{\rank}{\operatorname{rank}}
\newcommand{\calo}{\mathcal{o}}
\newcommand{\frako}{\mathfrak{o}}
\newcommand{\Cl}{\operatorname{Cl}}

\newcommand{\idim}{\operatorname{idim}}
\newcommand{\pdim}{\operatorname{pdim}}
\newcommand{\Ext}{\operatorname{Ext}}
\newcommand{\co}{\operatorname{co}}
\newcommand{\bfO}{\mathbf{O}}
\newcommand{\bfF}{\mathbf{F}} % Fitting Subgroup
\newcommand{\Syl}{\operatorname{Syl}}
\newcommand{\nor}{\vartriangleleft}
\newcommand{\noreq}{\trianglelefteqslant}
\newcommand{\subnor}{\nor\!\nor}
\newcommand{\Soc}{\operatorname{Soc}}
\newcommand{\core}{\operatorname{core}}
\newcommand{\Sd}{\operatorname{Sd}}
\newcommand{\mesh}{\operatorname{mesh}}
\newcommand{\sminus}{\setminus}
\newcommand{\diam}{\operatorname{diam}}
\newcommand{\Ass}{\operatorname{Ass}}
\newcommand{\projdim}{\operatorname{proj~dim}}
\newcommand{\injdim}{\operatorname{inj~dim}}
\newcommand{\gldim}{\operatorname{gl~dim}}
\newcommand{\embdim}{\operatorname{emb~dim}}
\newcommand{\hght}{\operatorname{ht}}
\newcommand{\depth}{\operatorname{depth}}
\newcommand{\ul}[1]{\underline{#1}}
\newcommand{\type}{\operatorname{type}}
\newcommand{\CM}{\operatorname{CM}}
\newcommand{\cech}[1]{\mathbin{\check{#1}}}
\newcommand{\cdim}{\operatorname{cdim}}
\newcommand{\Der}{\operatorname{Der}}
\newcommand{\trdeg}{\operatorname{trdeg}}
\newcommand{\calE}{\mathcal{E}}
\newcommand{\calF}{\mathcal{F}}
\newcommand{\calT}{\mathcal{T}}
\newcommand{\calN}{\mathcal{N}}
\newcommand{\calL}{\mathcal{L}}
\renewcommand{\Re}{\operatorname{Re}}
\renewcommand{\Im}{\operatorname{Im}}
\newcommand{\gr}{\operatorname{gr}}
\newcommand{\dist}{\operatorname{dist}}
\newcommand{\calH}{\mathcal{H}}


\geometry {
    margin = 1in
}

\titleformat
{\section}
[block]
{\Large\bfseries\sffamily}
{\S\thesection}
{0.5em}
{\centering}
[]


\titleformat
{\subsection}
[block]
{\normalfont\bfseries\sffamily}
{\S\S}
{0.5em}
{\centering}
[]


\begin{document}
\maketitle

In all the solutions, $\K$ denotes either $\R$ or $\bbC$. We shall work over $\K$ whenever possible and specialize to either field only when mentioned explicitly.

\section{Problem 1}

\subsection{Part (a)}

This is not true for $\bbC$-linear Hilbert spaces, indeed, let $T\colon\bbC\to\bbC$ denote the complex conjugation. Then clearly $T$ is an isometry but is not linear. Henceforth, I assume $\calH$ is an $\R$-linear Hilbert space and $T(0) = 0$. Note that without the second condition, any translation would be an ``isometry''. We have that $\|Tx - Ty\| = \|x - y\|$ and $T(0) = 0$, so that 
\begin{equation*}
    \|Tx\| = \|Tx - T(0)\| = \|x - 0\| = \|x\|\qquad\text{for all }~x\in\calH.
\end{equation*}

For $x, y\in\calH$, we have 
\begin{align*}
    \|x\|^2 - 2\langle x, y\rangle + \|y\|^2 &= \|x -  y\|^2\\
    &= \|Tx - Ty\|^2 \\
    &= \|Tx\|^2 - 2\langle Tx, Ty\rangle + \|Ty\|^2\\
    &= \|x\|^2 - 2\langle Tx, Ty\rangle + \|y\|^2,
\end{align*}
so that $\langle Tx, Ty\rangle = \langle x, y\rangle$ for all $x, y\in\calH$. Using this, for $x, y\in\calH$ and $\lambda\in\R$, we have 
\begin{align*}
    &~~~\|T(x + \lambda y) - Tx - \lambda Ty\|^2 \\
    &= \langle T(x + \lambda y) - Tx - \lambda Ty, T(x + \lambda y) - Tx - \lambda Ty\rangle\\
    &= \|T(x + \lambda y)\|^2 + \|Tx\|^2 + \|\lambda Ty\|^2 - 2\langle T(x + \lambda y), Tx\rangle - 2\langle T(x + \lambda y), \lambda T(y)\rangle + 2\langle Tx, \lambda Ty\rangle\\
    &= \|T(x + \lambda y)\|^2 + \|Tx\|^2 + \lambda^2\|Ty\|^2 - 2\langle T(x + \lambda y), Tx\rangle - 2\lambda\langle T(x + \lambda y), T(y)\rangle + 2\lambda\langle Tx, Ty\rangle\\
    &= \|x + \lambda y\|^2 + \|x\|^2 + \lambda^2\|y\|^2 - 2\langle x + \lambda y, x\rangle - 2\lambda\langle x + \lambda y, y\rangle + 2\lambda\langle x, y\rangle\\
    &= \|x + \lambda y - (x + \lambda y)\|^2 = 0,
\end{align*}
so that $T(x + \lambda y) = Tx + \lambda Ty$. Thus $T$ is a linear map from $\calH$ to $\calH$. Finally, since $\|Tx\| = \|x\|$ for all $x\in\calH$, $T$ is a bounded linear map with $\|T\| = 1$, as desired. 

\subsection{Part (b)}

Consider the map $F\colon\R\to\R^2$ given by $F(t) = (t, |t|)$. Note that $F(0) = (0, 0)$, and for $s, t\in\R$, we have 
\begin{equation*}
    \|F(t) - F(s)\|_{\R^2} = \left\|\left(t - s, |t| - |s|\right)\right\|_{\R^2} = \max\left\{|t - s|, \left||t| - |s|\right|\right\} = |t - s| = \|t - s\|_{\R},
\end{equation*}
where the third inequality follows from the fact that $|t - s|\ge\left||t| - |s|\right|$. This shows that $F$ is an isometry. But $F$ is not additive, since 
\begin{equation*}
    F(1) + F(-1) = (1, 1) + (-1, 1) = (0, 2)\ne (0, 0) = F(0).
\end{equation*}
Thus $F$ is an isometry but not a linear map.

\section{Problem 2}

Let $\K\in\{\R,\bbC\}$ and suppose $X$ is an NLS over $\K$. Fix $x\in X$. If $x\in Y$, then for all $\varphi\in X^\ast$ such that $\varphi|_Y = 0$, we have $\varphi(x) = 0$, so that $m(x) = 0 = \dist(x, Y)$. Suppose now that $x\notin Y$. If $\varphi\in X^\ast$ with $\|\varphi\|\le 1$ and $\varphi|_Y = 0$, then for all $y\in Y$, using the fact that $\varphi(y) = 0$, we have the inequality: 
\begin{equation*}
    |\varphi(x)| = |\varphi(x - y) + \varphi(y)| = |\varphi(x - y)|\le \|x - y\|,
\end{equation*}
so that taking infimums over $y\in Y$, we have 
\begin{equation*}
    |\varphi(x)|\le\inf_{y\in Y} \|x - y\| = \dist(x, Y).
\end{equation*}
This shows that $m(x)\le \dist(x, Y)$.

Define $\varphi\colon\K x + Y\to\K$ as $\varphi(\lambda x + y) = \lambda\dist(x, Y)$. This is clearly a linear map -- to see that this is bounded, note that for $\lambda\ne 0$, 
\begin{equation*}
    \left|\varphi(\lambda x + y)\right| = |\lambda|\dist(x, Y)\le |\lambda|\left\|x + \frac{1}{\lambda}y\right\|\le\|\lambda x + y\|,
\end{equation*}
and the above inequality clearly holds when $\lambda = 0$, since the left hand side is zero in that case. Therefore, $\|\varphi\|\le 1$ and is a bounded linear functional. Due to the Hahn-Banach Theorem, there is a norm-preserving extension of $\varphi$ to $X$, say $\wt\varphi\in X^\ast$. Then 
\begin{equation*}
    |\wt\varphi(x)| = \dist(x, Y)\implies m(x)\ge\dist(x, Y),
\end{equation*}
thereby implying the desired conclusion.

\section{Problem 3}

\subsection{Part (a)}

For each $x\in X\setminus\{0\}$, let $T_x\colon Y\to Z$ denote the bounded linear functional given by 
\begin{equation*}
    T_x(y) = \frac{f(x, y)}{\|x\|}.
\end{equation*}
For a fixed $y\in Y$, we have 
\begin{equation*}
    \sup_{x\in X\setminus\{0\}}\|T_x(y)\| = \sup_{x\in X\setminus\{0\}}\frac{\|f(x, y)\|}{\|x\|}.
\end{equation*}
Since $f(\cdot, y)$ is a bounded linear functional on $X$, there is a constant $C_y > 0$ such that $\|f(x, y)\|\le C_y\|x\|$ for all $x\in X$, so that 
\begin{equation*}
    \sup_{x\in X\sminus\{0\}}\|T_x(y)\|\le C_y.
\end{equation*}
Therefore, due to the Banach-Steinhaus Theorem, there is a constant $C > 0$ such that $\|T_x\|\le C$ for all $x\in X\setminus\{0\}$. In other words, for all $x\in X\setminus\{0\}$ and $y\in Y$, 
\begin{equation*}
    \frac{\|f(x, y)\|}{\|x\|}\le C\|y\|\implies \|f(x, y)\|\le C\|x\|\|y\|.
\end{equation*}
Finally, since $f(0, y) = 0$ for all $y\in Y$, we note that the above inequality holds for all $x\in X$ and $y\in Y$, as desired. 

\subsection{Part (b)}

Let $(x_0, y_0)\in X\times Y$, and let $\varepsilon > 0$. For $(x, y)\in X\times Y$, we have 
\begin{align*}
    \|f(x, y) - f(x_0, y_0)\| &= \left\|f(x, y) - f(x_0, y) + f(x_0, y) - f(x_0, y_0)\right\|\\ 
    &= \left\|f(x - x_0, y) + f(x_0, y - y_0)\right\|\\
    &\le \|f(x - x_0, y)\| + \|f(x_0, y - y_0)\|\\
    &\le C\|x - x_0\|\|y\| + C\|x_0\|\|y - y_0\|.
\end{align*}
Let $M > 0$ be a positive real number such that $\|x_0\|\le M$ and $\|y_0\|\le M$. Then, for $\|x - x_0\|\le 1$ and $\|y - y_0\|\le 1$, we have $\|x\|\le M + 1$ and $\|y\|\le M + 1$, so that 
\begin{equation*}
    \left\|f(x, y) - f(x_0, y_0)\right\|\le C(M + 1)\left(\|x - x_0\| + \|y - y_0\|\right).
\end{equation*}
Choose $\delta > 0$ such that 
\begin{equation*}
    \delta < \min\left\{1, \frac{\varepsilon}{2C(M + 1)}\right\}.
\end{equation*}
Then, for $\|x - x_0\| < \delta$ and $\|y - y_0\| < \delta$, we have $\|x\|\le M + \delta\le M + 1$ and $\|y\|\le M + \delta\le M + 1$, so that 
\begin{equation*}
    \|f(x, y) - f(x_0, y_0)\| < C(M + 1)\left(\frac{\varepsilon}{2C(M + 1)} + \frac{\varepsilon}{2C(M + 1)}\right) < \varepsilon.
\end{equation*}
Since $B_X(x_0, \delta)\times B_Y(y_0, \delta)$ forms a basic open set in $X\times Y$, it follows that $f\colon X\times Y\to Z$ is a continuous function.

\section{Problem 4}

To show that $A$ is continuous, we shall use the Closed Graph Theorem. To this end, we only need to show that $\operatorname{Graph}(A) = \left\{(x, Ax)\colon x\in\calH\right\}$ is closed in $\calH\times\calH$. Since $\calH\times\calH$ is a metric space, it suffices to show that every sequence in $\operatorname{Graph}(A)$ that converges in $\calH\times\calH$ must converge in $\operatorname{Graph}(A)$.

Let $(x_n, Ax_n)$ be a sequence in $\operatorname{Graph}(A)$ which converges to $(x, y)\in\calH\times\calH$. Since the projection operators $\operatorname{pr}_1,\operatorname{pr}_2\colon\calH\times\calH\to\calH$ are continuous, it follows that $x_n\to x$ and $Ax_n\to y$ in $\calH$. We shall show that $y = Ax$. Indeed, for any $z\in\calH$, using the fact that the linear functional $\langle\cdot, z\rangle\colon\calH\to\K$ is continuous,
\begin{equation*}
    \langle y, z\rangle = \lim_{n\to\infty}\langle Ax_n, z\rangle = \lim_{n\to\infty}\langle x_n, Az\rangle = \langle x, Az\rangle = \langle Ax, z\rangle.
\end{equation*}
That is, 
\begin{equation*}
    \langle Ax - y, z\rangle = 0\qquad\text{ for all } z\in\calH.
\end{equation*}
Set $z = Ax - y$ to obtain $\|Ax - y\| = 0$, so that $Ax = y$, as desired. Thus $\operatorname{Graph}(A)$ is closed in $\calH\times\calH$, and hence $A$ is continuous.

\section{Problem 5}

For each $y\in\calH$, consider the linear functional $\psi_x = B(\cdot, y)\colon\calH\to\K$. This is a bounded linear functional, since for all $x\in\calH$, 
\begin{equation*}
    \left|\psi_y(x)\right| = \left|B(x, y)\right|\le c\|x\|\|y\|,
\end{equation*}
so that $\|\psi_y\|\le c\|y\|$. Therefore, due to the Riesz Representation Theorem, there is a unique element $A(y)\in\calH$ such that $\langle x, Ay\rangle = B(x, y)$ for all $x\in\calH$. This defines a map $A\colon\calH\to\calH$. 

We contend that $A$ is a bounded linear functional. First to see that $A$ is linear, let $\alpha_1, \alpha_2\in\K$ and $y_1, y_2\in\calH$. Then for all $x\in\calH$, 
\begin{align*}
    \langle x, A(\alpha_1 y_1 + \alpha_2 y_2)\rangle &= B(x, \alpha_1y_1 + \alpha_2y_2)\\
    &= \overline\alpha_1 B(x, y_1) + \overline\alpha_2 B(x, y_2)\\
    &= \overline\alpha_1\langle x, Ay_1\rangle + \overline\alpha_2\langle x, Ay_2\rangle\\
    &= \langle x, \alpha_1 Ay_1\rangle + \langle x, \alpha_2 Ay_2\rangle\\
    &= \langle x, \alpha_1 Ay_1 + \alpha_2 Ay_2\rangle,
\end{align*}
so that $\langle x, \alpha_1 Ay_1 + \alpha_2 Ay_2 - A(\alpha_1 y_1 + \alpha_2 y_2)\rangle = 0$ for all $x\in\calH$. Setting $x = \alpha_1 Ay_1 + \alpha_2 Ay_2 - A(\alpha_1 y_1 + \alpha_2 y_2)$, it follows that 
\begin{equation*}
    A(\alpha_1y_1 + \alpha_2y_2) = \alpha_1 Ay_1 + \alpha_2 Ay_2.
\end{equation*}
To see that $A$ is bounded, note that for any $y\in\calH$, 
\begin{equation*}
    \|Ay\|^2 = \left|\langle Ay, Ay\rangle\right| = \left|B(Ay, y)\right|\le c\|y\|\|Ay\|.
\end{equation*}
If $Ay = 0$, then there is nothing to prove, but if $Ay\ne 0$, then $\|Ay\|\le c\|y\|$, whence $A$ is a bounded linear operator.

Now we show that the range $\mathcal R(A)$ of $A$ is closed. Indeed, suppose $(y_n)_{n = 1}^\infty$ is a sequence in $\calH$ such that $Ay_n\to z$ for some $z\in\calH$. We must show that $z\in\mathcal R(A)$. First, note that for any $u\in\calH$, 
\begin{equation*}
    b\|u\|^2\le\left|B(u, u)\right| = \left|\langle u, Au\rangle\right|\le \|u\|\|Au\|.
\end{equation*}
If $u\ne 0$, then dividing throughout by $\|u\|$, we obtain $\|Au\|\ge b\|u\|$. This inequality is trivially true if $u = 0$, and hence, it holds for all $u\in\calH$. This shows, in particular, that 
\begin{equation*}
    \|y_n - y_m\|\le \frac{1}{b}\|Ay_n - Ay_m\|.
\end{equation*}
Since $\left(Ay_n\right)_{n =1 }^\infty$ constitutes a convergent sequence in $\calH$, it is a Cauchy sequence, and in particular, for every $\varepsilon > 0$, there is a positive integer $N > 0$ such that whenever $m, n\ge N$, 
\begin{equation*}
    \|Ay_n - Ay_m\| < b\varepsilon\implies\|y_n - y_m\| < \varepsilon.
\end{equation*}
Hence, $(y_n)_{n = 1}^\infty$ forms a Cauchy sequence in $\calH$. Since $\calH$ is complete, there is a $y\in\calH$ with $y_n\to y$ as $n\to\infty$. Further, since $A$ is a continuous map, we must have that 
\begin{equation*}
    z = \lim_{n\to\infty} Ay_n = A\left(\lim_{n\to\infty}y_n\right) = Ay,
\end{equation*}
so that $z\in\mathcal R(A)$. This shows that $\mathcal R(A)$ is closed in $\calH$.

Since $\calH$ is a Hilbert space, there is a decomposition $\calH = \mathcal R(A) + \mathcal R(A)^\perp$. We shall show that $\mathcal R(A)^\perp = 0$, so that $\calH = \mathcal R(A)$. Let $\wt y\in\mathcal R(A)^\perp$. Then, 
\begin{equation*}
    0 = \left|\langle \wt y, A\wt y\rangle\right| = \left|B(\wt y, \wt y)\right|\ge b\|\wt y\|^2,
\end{equation*}
so that $\wt y = 0$, i.e., $\mathcal R(A)^\perp = 0$, equivalently, $\mathcal R(A) = \calH$, so that $A$ is surjective.

Finally, let $\varphi\in\calH^\ast$ be a bounded linear functional. Due to the Riesz Representation Theorem, there is a unique $u\in\calH$ such that for all $x\in\calH$, $\varphi(x) = \langle x, u\rangle$. Since $A$ is surjective, there is a $v\in\calH$ such that $u = Av$. Hence, for all $x\in\calH$, 
\begin{equation*}
    \varphi(x) = \langle x, u\rangle = \langle x, Av\rangle = B(x, v).
\end{equation*}
It remains to show that $v$ is unique. Suppose there is a $v'\in\calH$ with $\varphi(x) = B(x, v')$ for all $x\in\calH$. Then 
\begin{equation*}
    b\|v - v'\|^2\le \left|B(v - v', v - v')\right| = \left|B(v - v', v) - B(v - v', v')\right| = \left|\varphi(v - v') - \varphi(v - v')\right| = 0,
\end{equation*}
and hence, $\|v - v'\| = 0$, i.e., $v = v'$. This completes the proof.

\section{Problem 6}

Suppose $T$ has a right inverse, $S\colon Y\to X$ a bounded linear functional. Since $T\circ S = \id_Y$, for every $y\in Y$, we have 
\begin{equation*}
    T(S(y)) = \id_Y(y) = y,
\end{equation*}
so that $T$ is surjective. Further, for any $x\in X$, note that 
\begin{equation*}
    T\left(x - S(T(x))\right) = T(x) - T\left(S(T(x))\right) = T(x) - T(x) = 0,
\end{equation*}
so that $x - S(T(x))\in\calN(T)$, and we have a decomposition 
\begin{equation*}
    x = \left(x - S(T(x))\right) + S(T(x)).
\end{equation*}
Thus $X = \mathcal R(S) + \calN(T)$. Suppose now that $x\in\mathcal R(S)\cap\calN(T)$. Then there exists $y\in Y$ such that $x = S(y)$. But then $y = T(S(y)) = T(x) = 0$, so that $x = S(y) = 0$. Hence $\mathcal R(S)\cap\calN(T) = \{0\}$. 

It remains to show that $\mathcal R(S)$ is closed in $X$. Indeed, let $(x_n)_{n = 1}^\infty$ be a sequence in $\mathcal R(S)$ that converges to $x\in X$. We shall show that $x\in\mathcal R(S)$. For each $n\ge 1$, there exists a unique $y_n\in Y$ such that $x_n = S(y_n)$. Note that $y_n = T(S(y_n)) = T(x_n)$, and since $T$ is a continuous map, it must map convergent sequences to convergent sequences. Therefore, $y_n = T(x_n)\to y$ for some $y\in Y$. Moreover, since $S$ is continuous,
\begin{equation*}
    S(y) = S\left(\lim_{n\to\infty} y_n\right) = \lim_{n\to\infty} S(y_n) = \lim_{n\to\infty} x_n = x,
\end{equation*}
so that $x\in\mathcal R(S)$, i.e., $\mathcal R(S)$ is closed in $X$. This shows that $\calN(T)$ is complemented in $X$ by $\mathcal R(S)$.

Conversely, suppose $T$ is surjective and $\calN(T)$ is complemented in $X$ by a closed subspace $Z$. We shall show that $T$ has a right-inverse. Let $P\colon X\to Z$ denote the projection onto the closed subspace $Z$ and $Q\colon X\to\calN(T)$ that onto the closed subspace $\calN(T)$, which exist and is are bounded linear maps as we have seen in class. Let $\iota\colon Z\to X$ denote the inclusion. We contend that the map $U = T\circ\iota\colon Z\to Y$ is an isomorphism of Banach spaces. Being the composition of two bounded linear functionals, $U$ is a bounded linear functional. Further, if $U(z) = 0$, then $z\in\calN(T)$, but since $\calN(T)\cap Z = \{0\}$, it follows that $z = 0$. Hence, $U$ is injective. Further, since $T$ is surjective, for any $y\in Y$, there exists an $x\in X$ with $T(x) = y$. Let $z = P(x)$ and $w = Q(x)$. We have 
\begin{equation*}
    y = T(x) = T(P(x) + Q(x)) = T(P(x)) + T(Q(x)) = T(P(x)) = U(z).
\end{equation*}
Hence, $U$ is a bijective bounded linear map between Banach spaces, and due to the Open Mapping Theorem (or the Bounded Inverse Theorem), $U$ must be an isomorphism, and must admit a bounded linear inverse $U^{-1}\colon Y\to Z$. Set $S = \iota\circ U^{-1}$. We claim that $S$ is the desired right-inverse for $T$. Indeed, for any $y\in Y$, 
\begin{equation*}
    T(S(y)) = T(\iota(U^{-1}(y))).
\end{equation*}
Let $z\in Z$ be such that $U(z) = y$, therefore, $T(z) = y$, by construction. Hence, $T(S(y)) = T(z) = y$, as desired. This completes the proof.
\end{document}